\section{Data and method}
\label{sec:data}

This section states what each module is built from, what it holds fixed, and how a reader should
weigh what comes out of it. It supports every later section, and it sets the vocabulary used to
label results.

\paragraph{What the displacement engine is.} It is a conditional labor-income-loss stress test. A
share of the wage bill is removed by assumption, and the engine reports what follows for the
budget, for household credit and for individual balance sheets. It contains no model of
artificial intelligence, no adoption path and no forecast. Every displacement figure in this
paper is therefore conditional balance-sheet arithmetic, not a causal estimate and not a
prediction, and the level is an input rather than a result.

One level, ten percent of the wage bill, is reported throughout as the case with the least
extrapolation, and that phrase is meant literally rather than as a claim to observation. What
stays inside its observed range there is the reemployment relationship. The default uplifts are
transferred from a different episode, the allocation of displacement across households is assumed,
and the transmission from lost income to spending is a stated mapping. So the ten percent case is
a calibrated scenario whose weakest link is calibrated to data, not an observed or causally
identified result.

\subsection{Sources}

Claim levels and holder shares come from the Federal Reserve's financial accounts
\citep{FRBZ1}, read from the published data package. The instrument suffix in that release is not
the same across sectors, so bank, insurer and central bank holdings of one instrument sit under
different codes; matching on the exact suffix pushes whole holder classes into an unallocated
remainder. We match on the instrument family and take the closest suffix per sector, which cuts
the aggregate remainder from about a fifth to under two percent. Every series used is recorded
with the publisher's own description in the replication package.

Household balances come from the Survey of Income and Program Participation \citep{CensusSIPP},
household composition, earnings and occupation from the American Community Survey
\citep{CensusACS}, and official credit aggregates from the New York Federal Reserve
\citep{NYFedHHDC2026}. The wage share of adjusted gross income, and the realized gains bound that
travels with it, are from the Statistics of Income \citep{IRSSOI}; reemployment is fitted on
fourteen biennial vintages of the worker displacement release \citep{BLSWorkerDisplacement}; loss
rates and scenario severities are from the supervisory stress test \citep{FRBDFAST2026}; trust
fund income and depletion dates from the trustees' report \citep{SSATrustees2026}; the
distribution of corporate equity from the distributional financial accounts \citep{FRBDFA} and
behind them the Survey of Consumer Finances \citep{FRBSCF}; and the housing agency book from the
Enterprises' annual filings \citep{FannieMae10K2025,FreddieMac10K2025}. The AI side is built from
the filings of nine registrants, with the off-balance-sheet structures described by
\citet{BIS2026} absent from them. The wage bill and federal receipts used as bases are
\result{WageBillBn} and \result{FedReceiptsBn} billion dollars. Units are taken from the provider
at retrieval rather than assumed, and a discontinued series is replaced rather than carried.

\subsection{The household credit engine}

The engine maps a displacement level onto household credit losses without naming a set of
displaced households. Displacement is a probability applied to every household in the working
core, defined as a household with at least one employed member aged 25 to 64. The share of
earners displaced is set equal to the share of the wage bill displaced, and each household's
earners are drawn against that share.

A household that loses an earner does not default; it becomes more likely to default. The
uplift is five percentage points for one displaced earner and eight for two
\citep{GerardiHerkenhoffOhanianWillen2018}, and losses are exposure at default times loss given
default, with severities of 0.25 to 0.40 on first-lien mortgages, 0.80 to 1.00 on cards, 0.45 to
0.65 on auto and 0.75 to 1.00 on student debt. Applying a portfolio loss rate to exposure at
default instead would count the probability of default twice.

Survey balances are benchmarked to official aggregates first, and two ratios do different work.
The \emph{under-reporting factor} is the official aggregate over the full survey universe and is a
survey correction; the \emph{coverage share} is the working-core balance over that universe and is
a fact about which households the analysis covers. Table~\ref{tab:underreporting} gives both. The
identity between them holds to three decimals on every book, and all four factors were rebuilt
independently.

\begin{table}[htbp]
\centering
\caption{Survey balances benchmarked to official aggregates, 2026:Q2. The under-reporting factor
corrects the survey; the coverage share describes the analysis sample. Rebuilt.}
\label{tab:underreporting}
\setlength{\tabcolsep}{4pt}
\begin{tabular}{lrrrr}
\toprule
Book & Official aggregate (USD bn) & Survey universe (USD bn) & Under-reporting factor & Coverage share \\
\midrule
Mortgage & 13,100.0 & 10,403.3 & 1.2592 & 0.8205 \\
Credit card & 1,357.2 & 539.5 & 2.5159 & 0.7400 \\
Auto & 1,713.0 & 899.9 & 1.9036 & 0.7737 \\
Student & 1,650.0 & 1,159.3 & 1.4233 & 0.8832 \\
\bottomrule
\end{tabular}

\end{table}

One scope statement travels with every credit number here. The engine is first round only: house
prices, consumer spending, business revenue and the employment of everyone not displaced are held
fixed. Against a supervisory scenario, which moves all of those together, the figures are a lower
bound rather than an estimate.

\subsection{The fiscal module and the second round}

The fiscal module converts displacement into a dollar amount against the wage bill and applies a
loss per dollar with three parts: the labor tax no longer collected on wages that are not
re-earned, less the capital tax collected on the income that displaced them, plus any outlay
response. The share of wages re-earned is governed by a reemployment rate fitted on the
displacement vintages against labor market slack and solved as a fixed point, since displacement
changes the slack the fit is a function of. Losses are reported for the terminal year rather than
as a cumulative total divided by a horizon, and they are not linear in displacement: the share
re-earned falls as displacement rises. Trust fund effects use each fund's own payroll share of its
own income.

The second round is the weakest module evidentially and is labeled scenario wherever it appears.
The net fall in wage income becomes a demand shortfall through the difference between the
propensity to consume out of labor and out of capital income
\citep{MianStraubSufi2020,FagerengHolmNatvik2021}; the shortfall is expressed as a share of output
and compared with the fall in output in the supervisory severely adverse scenario; and the
supervisor's own published loss totals are scaled by that ratio. House prices enter through an
income elasticity \citep{HarterDreiman2004}, and employment through an Okun coefficient
\citep{BallLeighLoungani2017}. Three elements have no source: the shape of the loss curve beyond
the supervisory point, for which we carry a linear, a capped and a convex mapping; the propensity
ranges, measured on transitory shocks where displacement is persistent; and the path itself. We
report the whole grid of 648 combinations as a band and never a point inside it.

\subsection{Institution-level data}

The second round is applied across individual balance sheets rather than to an aggregate. The
panel is every insured bank and credit union at 30 June 2026, with the loan book by category and
the capital measure each is tested against: the tier 1 leverage ratio against its four percent
threshold for banks, the net worth ratio against its six percent threshold for credit unions.
After exclusions it is \result{NInstitutions} institutions, \result{NBanks} banks and
\result{NCreditUnions} credit unions, holding \result{SystemAssetsBn} billion dollars of assets.

\result{ExcludedInstitutions} institutions holding \result{ExcludedAssetsBn} billion dollars are
excluded because they report neither tier 1 capital nor total equity, being overwhelmingly insured
US branches of foreign chartered institutions whose capital sits at the parent; a capital test
shows them breaching before any loss, which is an artifact of a missing field. A capital test is
only informative if almost nothing fails it before any loss, and after the exclusion
\result{BaselineBreaches} institutions breach at baseline, \result{BaselinePctAssets} percent of
system assets. Losses are allocated pro rata by category holdings, so dispersion reflects business
mix alone and is a lower bound on true dispersion.

\subsection{Status labels, and what was independently rebuilt}
\label{sec:rebuild}

Three labels are used throughout and appear in every table caption. \emph{Rebuilt} means the
quantity was reproduced by an independent computational rebuild: a separate run that had no access
to this project's code and worked from raw data and a published specification alone, inside a
tolerance fixed before the rebuild began. \emph{Measured} means produced by the pipeline, clearing
every applicable bounds check, but not independently rebuilt. \emph{Scenario} means arithmetic
conditional on an assumption stated where the number appears. In prose we give the label where a
reader could otherwise mistake a scenario for a measurement; elsewhere the table caption carries
it.

What a rebuild establishes needs saying plainly, because the word invites more than it should. It
establishes computational reproducibility: that the stated inputs and the stated construction
produce the stated number. It says nothing about whether the construction is the right one, or
whether the economic assumptions behind it hold. A quantity can be rebuilt to six decimal places
and still rest on an assumption we have no evidence for, and several in this paper do.

Three rebuild rounds have been run. Round two scored \result{RepTwoScored} quantities,
\result{RepTwoSealed} of them against expected values fixed in advance;
\result{RepTwoFivePct} came within five percent and \result{RepTwoTolerance} inside the stated
tolerance, and \result{RepTwoMismatch} did not match. Round three covered the two constructions
the earlier rounds had not reached, the debt-only ratio and the assembled capital tax rate, and
matched all \result{RepThreeSealed} quantities, \result{RepThreeRounding} of them within
rounding. The blind was procedural rather than enforced: the rebuild ran on the same machine, with
the file of expected values reachable at a path the specification names, and reported it unopened
until the rebuild was final. The claim is therefore that these quantities were independently
rebuilt under a reported blind protocol, not that they were independently verified.
Appendix~\ref{app:rebuild} gives the round-by-round detail and the causes of the mismatches.

Separately, a standing bounds check states the bound a quantity must respect before it is
computed and tests it afterwards, running \result{PlausChecks} checks with
\result{PlausViolations} violations, all of them superseded rows retained deliberately.

\subsection{Disclosure of AI assistance}
\label{sec:aidisclosure}

The conceptual model, the definitions, the inclusion and exclusion rules, the equations, the
scenario definitions, the statistical choices and the interpretation of results are the authors'.
Generative AI tools provided computational assistance, in writing and checking analysis code, and
drafting assistance with the text; the independent computational rebuilds were carried out by
separate instances of such tools working from published specifications. The authors reviewed every
output and take responsibility for the content. No reported value originates in a model's
assertion: every number in the text and tables is generated from a measured artifact in the
replication package.
